%============================================================
% chiri_fill.tex  (main2.tex風デザイン + 穴埋め切替)
% 目的：claude_chiri.tex の内容を変えずに、main2.tex同様の
%       \answertrue/\answerfalse 切替で穴埋めプリント化する
% コンパイル例：
%   platex chiri_fill.tex
%   dvipdfmx chiri_fill.dvi
%============================================================
\documentclass[10pt,a4paper,dvipdfmx]{jsarticle}

\usepackage[utf8]{inputenc}
\usepackage[top=15mm, bottom=15mm, left=15mm, right=15mm]{geometry}
\usepackage{graphicx}
\usepackage{xcolor}
\usepackage{fancyhdr}
\usepackage{titlesec}
\usepackage{calc}
\usepackage{ulem}
\usepackage{tcolorbox}
\usepackage{pgffor}
\usepackage{otf}
\usepackage{multicol}
\usepackage{booktabs}
\usepackage{longtable}
\usepackage{array}
\usepackage{colortbl}
\usepackage{multirow}
\usepackage{amsmath}
\usepackage{enumitem}
\usepackage{pgfplots}
\usepackage{tikz}
\usetikzlibrary{patterns,arrows.meta,shapes.geometric}

\tcbuselibrary{skins,breakable}
\pgfplotsset{compat=1.18}

% ==============================================
% ★ 穴埋めモードの切り替え設定 ★
% 解答プリント（赤字）にする場合は \answertrue
% 生徒用プリント（空欄）にする場合は \answerfalse
\newif\ifanswer
\answertrue % ←ここを変更
% ==============================================

\newcounter{blanknum}

% 穴埋め用マクロ（番号＋下線で文字幅維持）
\newcommand{\aw}[1]{%
  \stepcounter{blanknum}%
  % 番号を左にはみ出させつつ、そのぶんだけ左余白を確保
  \ifnum\value{blanknum}<10
    \hspace{1.2em}%
  \else
    \ifnum\value{blanknum}<100
      \hspace{1.6em}%
    \else
      \hspace{1.9em}%
    \fi
  \fi
  \llap{\textsubscript{\scriptsize（\theblanknum）}}%
  \ifanswer
    \uline{\textcolor{red}{\textbf{\Large{#1}}}}%
  \else
    \uline{\textcolor{white}{\textbf{\Large{#1}}}}%
  \fi
}

%================================
% 色（claude_chiri.texの定義を維持）
%================================
\definecolor{mainblue}{RGB}{30,90,160}
\definecolor{maingreen}{RGB}{20,120,60}
\definecolor{mainorange}{RGB}{200,80,20}
\definecolor{mainred}{RGB}{180,30,30}
\definecolor{lightyellow}{RGB}{255,252,220}
\definecolor{lightblue}{RGB}{220,235,255}
\definecolor{lightgreen}{RGB}{220,250,230}
\definecolor{lightred}{RGB}{255,225,225}
\definecolor{lightgray}{RGB}{245,245,245}
\definecolor{darkgray}{RGB}{80,80,80}
\definecolor{tabhead}{RGB}{50,100,170}

%================================
% tcolorbox（claude_chiri.texの意匠を維持）
%================================
\tcbset{
  sectionbox/.style={
    enhanced, breakable,
    colback=lightblue, colframe=mainblue,
    fonttitle=\bfseries,
    title style={fill=mainblue},
    coltitle=white,
    top=2mm, bottom=2mm, left=3mm, right=3mm,
  },
  pointbox/.style={
    enhanced, breakable,
    colback=lightyellow, colframe=mainorange,
    fonttitle=\bfseries,
    title style={fill=mainorange},
    coltitle=white,
    top=2mm, bottom=2mm, left=3mm, right=3mm,
  },
  testbox/.style={
    enhanced, breakable,
    colback=lightgreen, colframe=maingreen,
    fonttitle=\bfseries,
    title style={fill=maingreen},
    coltitle=white,
    top=2mm, bottom=2mm, left=3mm, right=3mm,
  },
  warnbox/.style={
    enhanced, breakable,
    colback=lightred, colframe=mainred,
    fonttitle=\bfseries,
    title style={fill=mainred},
    coltitle=white,
    top=2mm, bottom=2mm, left=3mm, right=3mm,
  },
  notebox/.style={
    enhanced, breakable,
    colback=lightgray, colframe=darkgray,
    fonttitle=\bfseries,
    title style={fill=darkgray},
    coltitle=white,
    top=2mm, bottom=2mm, left=3mm, right=3mm,
  },
}

%================================
% 見出し（main2風：jsarticle+titlesec）
%================================
\titleformat{\section}[hang]
  {\Large\bfseries}{\colorbox{darkgray}{\color{white}\thesection}}{1em}{\uline}
\titleformat{\subsection}[hang]
  {\large\bfseries}{\colorbox{gray!30}{\thesubsection}}{1em}{}
\titleformat{\subsubsection}[hang]
  {\normalsize\bfseries}{}{0em}{■ }

%================================
% ヘッダー・フッター（main2風）
%================================
\pagestyle{fancy}
\fancyhf{}
\fancyhead[L]{\textbf{地理 中間試験対策まとめ}}
\fancyhead[R]{\textbf{林業・水産業・鉱産資源}}
\fancyfoot[C]{\thepage}
\renewcommand{\headrulewidth}{0.4pt}

%================================
% 便利コマンド（内容は変えず、\key/\imp を穴埋め化）
%================================
\newcommand{\rank}[1]{\textcircled{\scriptsize #1}}
\newcommand{\ans}[1]{\underline{\hspace{#1}}}

% 方針：\key と \imp の両方を穴埋めにする（本文は変更しない）
\newcommand{\key}[1]{\aw{#1}}
\newcommand{\imp}[1]{\aw{#1}}

%================================
% 表スタイル（claude_chiri.tex互換）
%================================
\newcolumntype{C}[1]{>{\centering\arraybackslash}p{#1}}
\newcolumntype{L}[1]{>{\raggedright\arraybackslash}p{#1}}
\newcolumntype{R}[1]{>{\raggedleft\arraybackslash}p{#1}}

\begin{document}

%============================================================
% タイトル（内容維持、見た目をmain2寄せ）
%============================================================
\begin{center}
  {\Huge \textbf{地理 中間試験対策まとめプリント}}\\[0.5em]
  {\Large 林業・水産業・鉱産資源・資源エネルギー}\\[0.3em]
  {\small 〔範囲〕グラフ読解 / 共通テスト過去問類題 / 自由記述対策 / 統計判定ポイント}
\end{center}
\vspace{1em}

%============================================================
% ここから本文：claude_chiri.tex をそのまま貼り付け（内容不変）
% ※ このコミットでは冒頭のみを移植しています。
%    続きも全て移す場合は、claude_chiri.tex の本文全体を
%    \section{林業} 以降〜\end{document}直前まで貼り付けてください。
%============================================================

%%============================================================
\section{林業}
%%============================================================

\subsection{基礎知識：用語整理}

\begin{tcolorbox}[sectionbox, title={■ 林業の基本用語}]
\begin{description}[leftmargin=2.5cm, style=nextline, labelsep=4mm]
  \item[\key{林業}] 豊かな森林資源を活用し，木製品の原材料（用材），燃料（薪炭材），
    きのこ・天然ゴム・樹皮など非木材林産物を生産する産業。
  \item[\key{用材}] 家具・建築材料になる丸太，製材品，合板，パルプ等。
    産業用丸太 ≈ 約20億m³/年（2022）。
  \item[\key{薪炭材}] 薪や炭などの燃料用丸太。約20億m³/年（2022）。
    \imp{発展途上国}で全木材生産の6割以上を占める。
  \item[\key{丸太}（用材広義）] 産業用丸太（用材）＋薪炭材を合わせた概念。
  \item[\key{合板}] 薄くスライスした木の単板を何層も積層・重ね合わせたもの。
    加工しやすく強度・断熱性・吸音性に優れる。型枠にも利用。
  \item[\key{パルプ}] 木材チップを薬品で溶かした繊維分。紙の原料。草・わら・さとうきびからも製造可。
  \item[\key{チーク材}] 熱帯の常緑広葉樹。高級家具・内装用。東南アジア産。
  \item[\key{ラワン材}] 熱帯産の柔らかい広葉樹。合板に加工，コンクリート型枠に使用。
  \item[\key{人工林}] 経済発展で天然林が伐採・需要増→人工林が増加。日本国土の約4割。
  \item[\key{退耕還林}] 中国の政策：農地を森林に戻す政策。\key{生態移民}と連動。
  \item[\key{生物多様性}] 森林は陸域生物種の約8割が生息。保全が重要課題。
\end{description}
\end{tcolorbox}

\vspace{2mm}

\subsection{世界の森林分布と樹種}

\begin{tcolorbox}[pointbox, title={■ 主要な森林タイプと分布}]
\begin{tabular}{L{3cm}L{4.5cm}L{5cm}}
\toprule
\rowcolor{mainorange!20}
\textbf{森林タイプ} & \textbf{主な分布地域} & \textbf{特徴・用途} \\
\midrule
\key{熱帯雨林}（常緑広葉樹） & アマゾン，コンゴ盆地，東南アジア & チーク・ラワン；二酸化炭素吸収量大 \\
\addlinespace
\key{タイガ}（亜寒帯針葉樹林） & カナダ，ロシア，北欧 & 建築材・パルプ原料；針葉樹林 \\
\addlinespace
\key{温帯林}（温帯広葉樹・混交林） & 西・中央ヨーロッパ，日本，中国東部 & 広葉樹・針葉樹が混在 \\
\addlinespace
\key{硬葉樹林}（地中海性） & 地中海沿岸，カリフォルニア & コルク樫など；乾燥耐性 \\
\bottomrule
\end{tabular}

\medskip
\begin{itemize}[leftmargin=1em]
  \item 世界の森林面積は\imp{40億ha}超（陸地面積の約3割）。9割以上が天然林。
  \item 亜寒帯林（タイガ）は建築資材・パルプの原料として重要。
  \item 熱帯雨林の伐採→農地転用→\imp{二酸化炭素吸収量の減少}→地球温暖化の要因。
\end{itemize}
\end{tcolorbox}

\end{document}
